% Laget automatisk av elevark/lag_elevark.py fra opplegg/kraftsensor.html. Endringer her blir overskrevet.
\documentclass[11pt]{article}
\usepackage{elevark}
\begin{document}
\elevarkhode{Fysikk 1 · Bevegelsesmengde og støt}{Impuls med kraftsensor}{https://lektorbo.no/opplegg/kraftsensor.html}{Opplegget på nett}

En vogn kjører inn i en kraftsensor og spretter tilbake. Er arealet under grafen for kraft mot tid lik endringen i bevegelsesmengden til vogna?

\section*{Hensikt}

\begin{itemize}
\item Måle kraften i et kort støt og finne impulsen som arealet under F(t)-grafen.
\item Undersøke om impulsen er lik endringen i bevegelsesmengde: $F \cdot Δt = Δp$.
\item Se hvordan en mykere støtfanger endrer støttiden og den største kraften.
\end{itemize}

\section*{Utstyr}

\begin{itemize}
\item Dynamikkbane og vogn, for eksempel Smart Cart med innebygd kraftsensor, eller en vanlig vogn og en kraftsensor festet i enden av banen.
\item Fjærstøtfanger eller magnetstøtfanger. Gjerne også en myk støtfanger (skumgummi) til sammenligning.
\item Bevegelsessensor, eller fartsmåling fra Smart Cart.
\item Vekt og datamaskin med Capstone, SPARKvue eller et annet program som logger målingene.
\end{itemize}

\section*{Framgangsmåte}

\begin{enumerate}
\item Vei vogna. Juster banen slik at den står vannrett: Vogna skal bli stående i ro når du slipper den.
\item Still loggefrekvensen for kraftsensoren høyt, minst 500~Hz. Støtet varer bare noen titalls millisekunder.
\item Nullstill kraftsensoren mens ingenting berører den.
\item Gi vogna et forsiktig dytt mot sensoren, og logg både kraft og fart til vogna har sprettet tilbake.
\item Les av farten rett før og rett etter støtet. Husk fortegn.
\item Eksporter kraftmålingene som tekst (tid og kraft i to kolonner). Lim dem inn i verktøyet på nettsiden, eller bruk programmets eget verktøy for areal under grafen.
\item Gjenta med en mykere støtfanger og omtrent samme fart.
\end{enumerate}

\section*{Oppgaver}

\begin{oppgaveliste}
\oppgave \fase{gjett}{Før}Skisser hvordan du tror grafen for kraft mot tid ser ut for en fjærstøtfanger og for en myk støtfanger. Vogna har samme fart begge gangene.
\svar
\oppgave \fase{test}{Regn}Regn ut $Δp$ for vogna. Pass på fortegnet: Vogna snur.
\svar
\oppgave \fase{test}{Mål}Finn arealet under kraftgrafen i støtet. Hvor stort er avviket fra $Δp$ i prosent?
\svar
\oppgave \fase{forklar}{Forklar}Sammenlign støttiden og den største kraften for de to støtfangerne. Stemmer det med skissen din? Hvorfor er arealet omtrent det samme?
\svar
\oppgave \fase{forklar}{Vurder}Hvilke feilkilder kan gi avvik mellom impulsen og $Δp$? Tenk på loggefrekvens, nullstilling, friksjon og hvor nøyaktig farten leses av.
\svar
\oppgave \fase{utfordring}{Utfordring}Regn ut middelkraften i støtet som $Δp/Δt$. Hvorfor er den mindre enn den største kraften? Hvilken av dem er viktigst for om noe går i stykker?
\svar
\end{oppgaveliste}

\end{document}
